In this section we describe the category $\T$ of 
$\Uq=\Uq(\mathfrak{sl}_2)$-tilting modules for the quantum enveloping 
algebra\footnote{We like to work over the cyclotomic field $\bC$, but any field of characteristic zero containing $q$ would work.} of $\mathfrak{sl}_2$ at a fixed root of unity $q$.
We give a hopefully self-contained summary of the results in the $\mathfrak{sl}_2$ case 
since most results are either spread out over the 
literature or only mentioned implicitly. But 
we have also included some new observations related to our context.



We mostly follow~\cite[Chapters~1,~2,~3]{ja} with our 
notation, $v$ denotes an indeterminate, $q\in\bC$ denotes a 
fixed root of unity and all module categories 
in this section are categories of \textit{left} 
modules.




\subsection{Quantum groups at roots of unity}\label{sub-qroot}
Before we start, let us fix some notions.
Given $a\in\bZ$ and $b\in\bN$, let
\begin{gather*}
[a]=\frac{v^a-v^{-a}}{v-v^{-1}}=v^{a-1}+v^{a-3}+\dots+v^{-a+1}+v^{-a+1},\quad\quad[b]!=[1][2]\cdots [b-1][b],\\
\qbin{a}{b}=\frac{[a][a-1]\cdots[a-b+2][a-b+1]}{[b]!}\in\bZ[v,v^{-1}].
\end{gather*}
be the \textit{quantum integer}, the \textit{quantum factorial} and 
the \textit{quantum binomial}.
By convention, $[0]!=1$.


\begin{defn}\label{def-qv} The \textit{quantum special linear algebra} $\Uv(\mathfrak{sl}_2)$ 
is the associative, unital $\bQ(v)$-algebra 
generated by $K,K^{-1},E$ and $F$ subject to the relations
\begin{gather*}
KK^{-1}=K^{-1}K=1,\quad\quad KE=v^2EK,\quad\quad KF=v^{-2}FK,\quad\quad
EF-FE =\dfrac{K-K^{-1}}{v-v^{-1}}.
\end{gather*}
We write $\Uv=\Uv(\mathfrak{sl}_2)$ for short.
\end{defn}


Recall that $\Uv$ is a Hopf algebra with coproduct $\Delta$, 
antipode $S$ and the counit $\varepsilon$ given by
\begin{gather*}
\Delta(E)=E\otimes 1+K\otimes E,\quad\quad\Delta(F)=F\otimes K^{-1}+1\otimes F,\quad\quad\Delta(K)=K\otimes K,\\
S(E)=-K^{-1}E,\;\;\;\; S(F)=-FK,\;\;\;\; S(K)=K^{-1},\quad\quad\varepsilon(E)=\varepsilon(F)=0,\;\;\;\;\varepsilon(K)=1.
\end{gather*}
This allows to extend actions 
to tensor products and 
to duals, and there is a trivial $\Uv$-module.



We want to ``specialize'' the $v$ to be a root of unity $q\in\bC$. 
To this end, 
let $\mathrm{A}=\bZ[v,v^{-1}]$ and we consider 
Lusztig's $\mathrm{A}$-form $\Ua$ from~\cite{lu2}.


\begin{defn}(\textbf{Lusztig's $\mathrm{A}$-form $\Ua$})\label{def-qA} Define for 
$j\in\bZ_{>0}$ the \textit{$j$-th divided powers}
\[
E^{(j)}=\frac{E^{j}}{[j]!},\quad\quad F^{(j)}=\frac{F^{j}}{[j]!}.
\]
$\Ua$ is 
defined as the $\mathrm{A}$-subalgebra of 
$\Uv$ generated by $K,K^{-1}$, $E^{(j)}$ and $F^{(j)}$.
\end{defn}




\begin{defn}\label{def-qq} Fix a root of unity 
$q\in\bC$, $q\neq \pm 1$ and denote by 
$l$ the order\footnote{The square is only important for roots 
of unity of even order, but not 
for roots of unity of odd order (for these the order of $q^2$ is the same as 
the order of $q$), e.g. we have $l=3$ for third as well as for sixth roots of unity.} of $q^2$. Consider 
$\bC$ as an $\mathrm{A}$-module by specializing $v$ to $q$. Define
\[
\Uq=\Ua\otimes_{\mathrm{A}}\bC.
\]
We abuse notation and write $E^{(j)}$ instead 
of $E^{(j)}\otimes 1$. Analogously for the other generators.
\end{defn}




\begin{rem}\label{rem-indhopf}
It is easy to check that $\Ua$ is a Hopf 
subalgebra of $\Uv$. Thus, $\Uq$ inherits a Hopf algebra structure from $\Uv$. 
In particular, if one has a $\Uq$-module $M$, then 
$M^*=\Hom_{\Uq}(M,\bC)$ has the usual induced action. 
It follows that, if $m\in M$ is an eigenvector of $K$ with 
eigenvalue $\alpha$ and $m^*\in M^*$ a dual vector with respect to 
some $K$-stable complement of $\bC m$ in $M$, then
\begin{equation}\label{eq-eigen}
Km=\alpha m\Longleftrightarrow Km^*=\alpha^{-1}m^*.
\end{equation}
Moreover, note that $[j]=0\in\Uq$ iff $l|j$. 
This implies 
$E^l=[l]!E^{(l)}=0$ and $F^l=[l]!F^{(l)}=0$. 
It is also true that $K^{2l}=1$, see~\cite[Lemma~4.4~(a)]{lu1}.
\end{rem}




\subsection{Weyl modules, dual Weyl modules and simple modules}\label{sub-wsmod}



\begin{defn}(\textbf{Weyl, dual Weyl and simple modules})\label{defn-weyl}
Let $i\in\bN$ and denote 
by $\Vn i$ the $i$\textit{-th Weyl module}. 
This is the $\bC$-vector space 
with basis $m_0,\dots,m_i$ and an $\Uq$-action defined by
\[
Km_k=q^{i-2k}m_k,\quad\quad E^{(j)}m_k=\qbin{i-k+j}{j}m_{k-j},\quad\quad F^{(j)}m_k=\qbin{k+j}{j}m_{k+j},
\]
with the convention that $m_{<0}=m_{>i}=0$. 
The \textit{$i$-th dual Weyl module}, 
denoted by $\dVn i$, is obtained from 
$\Vn i$ by taking the dual, that is, 
$\Hom_{\mathbb{Q}(q)}(\Vn i,\bC)=(\Delta_q(i))^*=\dVn i$.
\end{defn}


It is easy to check directly that 
$\Vn i$ has a \textit{unique simple head} $\Ln i$ which 
is also the \textit{unique simple socle} of $\dVn i$. 
See also~\cite[Section~4]{apw}.



\begin{ex}\label{ex-weyl}
Let $q=\frac{-1+\sqrt{-3}}{2}$ (we use this $q$ in all examples with $l=3$ in what follows).



The Weyl module $\Vn 3$ can be visualized as
\begin{equation}\label{eq-example1}
\xymatrix{
  m_3  \ar@<2pt>[r]^{+1}\ar@(ul,ur)|{q^{-3}}  &  {\color{colormy}m_2}  \ar@[mycolor]@<2pt>[l]^{{\color{mycolor} 0}}\ar@[colormy]@<2pt>[r]^{{\color{colormy} -1}}\ar@[colormy]@(ul,ur)|{{\color{colormy}q^{-1}}} & {\color{colormy}m_1}  \ar@[mycolor]@<2pt>[r]^{{\color{mycolor} 0}}\ar@[colormy]@<2pt>[l]^{{\color{colormy}-1}}\ar@[colormy]@(ul,ur)|{{\color{colormy}q^{+1}}}  &  m_0.  \ar@<2pt>[l]^{+1}\ar@(ul,ur)|{q^{+3}} \\
}
\end{equation}
Here the action of 
$E$ points right, the action of 
$F$ left, $E^{(3)}m_3=m_0$, $F^{(3)}m_0=m_3$, $K$ acts as a loop 
and all other actions are zero. 
Thus, the $\mathrm{A}$-span of $\{m_1,m_2\}$ is now, 
in contrast to the classical case, stable under the action of $\Uq$. The complement 
however is not an $\Uq$-submodule.



Another example (where we have excluded the actions of the divided powers) is $\Vn{4}$:
\begin{equation}\label{eq-example2}
\xymatrix{
  m_4  \ar@<2pt>[r]^{+1}\ar@(ul,ur)|{q^{-4}}  &  m_3  \ar@<2pt>[l]^{+1}\ar@<2pt>[r]^{-1}\ar@(ul,ur)|{q^{-2}} & {\color{green} m_2}  \ar@[mycolor]@<2pt>[r]^{{\color{mycolor} 0}}\ar@[red]@<2pt>[l]^{{\color{mycolor} 0}}\ar@[colormy]@(ul,ur)|{{\color{colormy}q^{0}}} & m_1  \ar@<2pt>[r]^{+1}\ar@<2pt>[l]^{-1}\ar@(ul,ur)|{q^{+2}} &  m_0.  \ar@<2pt>[l]^{+1}\ar@(ul,ur)|{q^{+4}} \\
}
\end{equation}
We note that $\Vn{4}$ has the 
trivial $\Uq$-module spanned by $m_2$ as a $\Uq$-submodule.
\end{ex}




$K$ acts on $\Vn i$ via the eigenvalues $q^{-i},q^{-i+2},\dots,q^{i-2},q^{+i}$. 
Thus, by~\eqref{eq-eigen}, the same is 
true for $\dVn i$. Moreover, the 
$\Ln i$ are self-dual, see e.g.~\cite[Section~4]{apw}.




\begin{prop}\label{prop-weyl}
We have the following.
\begin{itemize}
\item[(a)] $\Vn i\cong\Ln i$ iff $i<l$ or $i\equiv -1\mod l$.
\item[(b)] Suppose $i=al+b$ for 
some $a,b\in\bN$ with $b\leq l-2$. 
Set $i^{\prime}=(a+2)l-b-2$. Then there exists a short exact sequence
\[
0\longrightarrow \Ln i\hooklongrightarrow\Vn{i^{\prime}}\twoheadlongrightarrow \Ln{i^{\prime}}\longrightarrow 0.
\]
Moreover, $\Ln{i^{\prime}}$ is the 
head and $\Ln i$ is the socle of $\Vn{i^{\prime}}$.\makeqed
\end{itemize}
\end{prop}


\begin{proof}
This is~\cite[Corollary~4.6]{apw}.
\end{proof} 
\begin{cor}\label{cor-weyl}
We have the following.
\begin{itemize}
\item[(a)] $\dVn i\cong\Ln i$ iff $i<l$ or $i\equiv -1\mod l$.
\item[(b)] Suppose $i=al+b$ for some 
$a,b\in\bN$ with $b\leq l-2$. Set $i^{\prime}=(a+2)l-b-2$. 
Then there exists a short exact sequence
\[
0\longrightarrow L_q(i^{\prime})\hooklongrightarrow\nabla_q(i^{\prime})\twoheadlongrightarrow \Ln i\longrightarrow 0.
\]
Moreover, $\Ln{i^{\prime}}$ is the socle and 
$\Ln i$ is the head of $\dVn{i^{\prime}}$.\makeqed
\end{itemize}
\end{cor}



\begin{proof}
This follows from $\Ln i\cong(\Ln i)^*$ and 
the fact that ${}^*$ is an exact, contravariant functor.
\end{proof}



\begin{ex}\label{ex-weyl2}
For $i=0$ and $l=3$ we have $i^{\prime}=4$. 
The trivial $\Uq$-module $\Ln 0$ appears as a 
submodule of $\Vn 4$, compare to~\eqref{eq-example2}.
\end{ex}



We should note here that 
there are two different 
types of $\Uq$-modules known as \textit{type} $1$ \textit{and} $-1$, 
see for example~\cite[Section 1]{apw}. 
For us the difference between the two types is not important 
in this paper and we \textit{only} consider $\Uq$-modules of type $1$. 
The (more general) treatment in~\cite[Section 1]{apw} 
ensures that the consideration of only type 
$1$ is still enough to get results for both types.



\begin{cor}\label{cor-simples}
The set $\{\Ln i\mid i\in\bN\}$ is a  
complete set of simple, pairwise non-isomorphic, 
finite-dimensional $\Uq$-modules (of type $1$).\makeqed
\end{cor}


\begin{proof}
See~\cite[Corollary~6.3]{apw}.
\end{proof}



\begin{rem}\label{rem-semisimple}
The category of $\Uq$-modules is 
\textit{far} from being semisimple: Proposition~\ref{prop-weyl} 
says that $\Vn i$ is never simple whenever we are in the case (b). 
See also~\eqref{eq-example1}.
\end{rem}



\subsection{Tilting modules and the tilting category \texorpdfstring{$\T$}{T}}\label{sub-tiltmod}
\begin{defn}(\textbf{$\Delta$- and $\nabla$-filtration})\label{defn-filt}
We say that a $\Uq$-module $M$ has 
a $\Delta$\textit{-filtration} if there exists a descending sequence of 
$\Uq$-submodules
\[
M=F_{0}\supset F_1\supset\cdots\supset F_i\supset\dots,\quad\quad\bigcap_{i=0}^{\infty}F_i=0,
\]
such that for all $i=0,1\dots$ we have 
$F_{i}/F_{i+1}\cong \Vn{i^{\prime}}$ for 
some $i^{\prime}\in\bN$. A $\nabla$\textit{-filtration} is 
defined similarly, but using 
$\dVn{i^{\prime}}$ instead of 
$\Vn{i^{\prime}}$ and an ascending sequence of $\Uq$-submodules, that is
\[
0=F_{0}\subset F_1\subset\cdots\subset F_i\subset\dots,\quad\quad\bigcup_{i=0}^{\infty}F_i=M,
\]
such that for all $i=0,1\dots$ we 
have $F_{i+1}/F_{i}\cong \dVn{i^{\prime}}$ for some $i^{\prime}\in\bN$.
\end{defn}


One can prove, following similar arguments as in~\cite[Proposition~II.4.16]{jarag},
that such filtrations are unique up to reordering. It is clear that 
a finite-dimensional $\Uq$-module $M$ has a 
$\Delta$-filtration iff $M^*$ has a $\nabla$-filtration.



\begin{defn}(\textbf{Tilting modules})\label{defn-tilt}
We call a $\Uq$-module $M$ a $\Uq$-\textit{tilting module} 
if it has a $\Delta$- and a $\nabla$-filtration. 
We say for short that $M$ is \textit{tilting}.
\end{defn}




\begin{ex}\label{ex-tilt}
Recall that the category of finite-dimensional $\Uv$-modules is 
semisimple. It follows that all finite-dimensional 
$\Uv$-modules are tilting. Moreover, for $i<l$ or $i\equiv -1 \mod l$, 
by Proposition~\ref{prop-weyl} and 
Corollary~\ref{cor-weyl}, we see that $\Vn i\cong\Ln i\cong\dVn i$ is tilting.
\end{ex}



It is easy to see 
(in our $\mathfrak{sl}_2$ case - in general this 
is non-trivial, see e.g.~\cite[Theorem~3.3]{par}) 
that $\Vn i\otimes_{\bC}\Vn{i^{\prime}}$ has a $\Delta$-filtration with factors 
(likewise for $\dVn i\otimes_{\bC}\dVn{i^{\prime}}$)
\begin{equation}\label{eq-tiltfact}
\Vn{|i-i^{\prime}|},\,\Vn{|i-i^{\prime}|+2},\,\dots, \Vn{i+i^{\prime}-2},\,\Vn{i+i^{\prime}},\quad\quad i,i^{\prime}\geq 0.
\end{equation}



\begin{defn}(\textbf{Full tilting category})\label{defn-ftiltcat}
We denote by $\T^{\mathrm{all}}$ the full 
subcategory of all $\Uq$-modules that are tilting. 
We denote its class of objects by $\Ob(\T^{\mathrm{all}})$.
\end{defn}



\begin{prop}\label{prop-basicp}
The category $\T^{\mathrm{all}}$ has the following properties.
\begin{itemize}
\item[(a)] If $M,M^{\prime}\in\Ob(\T^{\mathrm{all}})$, then $M\oplus M^{\prime}\in\Ob(\T^{\mathrm{all}})$.
\item[(b)] If $M\oplus M^{\prime}\in\Ob(\T^{\mathrm{all}})$, then $M,M^{\prime}\in\Ob(\T^{\mathrm{all}})$.
\item[(c)] If $M,M^{\prime}\in\Ob(\T^{\mathrm{all}})$, then $M\otimes_{\bC} M^{\prime}\in\Ob(\T^{\mathrm{all}})$.\makeqed
\end{itemize}
\end{prop}



\begin{proof}
Part (a) is immediate. For part (b) one can, for 
example, argue as in~\cite[Proposition~3.7]{hum}, and the non-trivial 
part (c) can be deduced from~\eqref{eq-tiltfact} 
and the corresponding statement for the dual Weyl modules $\dVn i$.
\end{proof}


\begin{defn}(\textbf{The tilting category $\T$})\label{defn-tiltcat}
Let $\T$ be the full subcategory of $\T^{\mathrm{all}}$ consisting of:
\begin{itemize}
\item The objects $\Ob(\T)$ are all 
finite-dimensional $\Uq$-tilting modules $M\in\Ob(\T^{\mathrm{all}})$.
\item The morphisms $\Hom_{\T}(M,N)$ are all $\Uq$-intertwiners $f\in\Hom_{\Uq}(M,N)$. 
\end{itemize}
The objects of $\T$ have finite filtrations, 
i.e. $F_N=F_{N+i}$ for some $N\in\bN$ and all $i\in\bN$.
\end{defn}

Recall that an additive category is \textit{Krull-Schmidt}, 
if each object can be uniquely decomposed (up to permutation) 
into a finite direct sum of \textit{indecomposable objects}, i.e. objects 
$O$ such that $O\cong O^{\prime}\oplus O^{\prime\prime}$ implies that 
$O^{\prime}\cong 0$ or $O^{\prime\prime}\cong 0$.

\begin{lem}\label{lem-tiltcat}
We have the following.
\begin{itemize}
\item[(a)] The category $\T$ is additive (but not abelian), closed under 
finite direct sums, finite tensor products and under duals.
\item[(b)] $\T$ is a Krull-Schmidt category whose 
indecomposable objects are parametrized by $\bN$.
\item[(c)] $\Vn i$, $\dVn i$ and $\Ln i$ are in $\T$ 
iff $i<l$ or $i\equiv-1\mod l$.\makeqed
\end{itemize}
\end{lem}


\begin{proof}
%Most of these statements can be verified by following 
%the arguments which prove similar statements 
%in related categories, see for example~\cite[Section~1]{hum}.
%
\textbf{(a).} From Proposition~\ref{prop-basicp} and the fact that 
$\Hom_{\Uq}(\cdot,\bC)$ commutes with finite sums.

\textbf{(b).} 
The Krull-Schmidt property follows from finite-dimensionality and 
(b) of Proposition~\ref{prop-basicp}.
By Proposition~\ref{prop-tilt} below, all finite-dimensional 
indecomposable tiltings are of the form $\Tn i$.


\textbf{(c).} A direct consequence of Proposition~\ref{prop-weyl} and Corollary~\ref{cor-weyl}.
\end{proof}



\begin{rem}\label{rem-ribbon}
It follows from Lemma~\ref{lem-tiltcat} part (a) 
that $\T$ is a \textit{rigid category} 
(a monoidal category with duals and certain compatibility properties). 
Moreover, $\T$ is even a \textit{ribbon (tensor) category} and 
$\T$ gives rise to a \textit{modular category} 
(roughly: one mods out by tiltings whose quantum trace is zero) 
and thus, gives a $2+1$-dimensional TQFT and can 
be used to define the \textit{Witten-Reshetikhin-Turaev invariants} of $3$-manifolds. 
Good treatments of this are~\cite[Section~4]{an1} or~\cite[Section~7]{saw}.
\end{rem}



Define, using Propositions~\ref{prop-weyl} 
and~\ref{prop-basicp} and~\eqref{eq-tiltfact}, 
a family $(\Tn i)_{i\in\bN}$ of indecomposable tiltings as follows. 
We start by setting $\Tn 0=\Vn 0$ and $\Tn 1=\Vn 1$. 
Then we denote by $m_q\in\Tn 1$ any eigenvector for $K$ with eigenvalue $q$. 
For each $i>1$ we define $\Tn i$ to be the indecomposable 
summand of $(\Tn 1)^{\otimes i}$ which contains 
the vector $m_q\otimes\cdots\otimes m_q\in (\Tn 1)^{\otimes i}$.



Denote by $(T:\Vn{i})\in\bN$ the \textit{filtration multiplicity} for an $\Uq$-tilting module $T$. It is clear, by using~\eqref{eq-tiltfact}, that
\begin{equation}\label{eq-tiltfact2}
((\Tn 1)^{\otimes i}:\Vn i)=1,\quad\quad((\Tn 1)^{\otimes i}:\Vn{i^{\prime}})=0,\quad\quad\text{for } i^{\prime}>i.
\end{equation}
Hence, $\Tn i$ may also be described as 
the unique indecomposable summand of 
$(\Tn 1)^{\otimes i}$ that contains $\Vn i$. We note the following more precise statement.



\begin{prop}\label{prop-tilt}
\item[(a)] $\Vn i\cong\Ln i\cong\Tn i\cong\dVn i$ iff $i<l$ or $i\equiv -1\mod l$.
\item[(b)] Suppose $i=al+b$ for some $a,b\in\bN$ with 
$b\leq l-2$. Set $i^{\prime}=(a+2)l-b-2$. Then there exist short exact sequences
\[
0\longrightarrow \Vn{i^{\prime}}\hooklongrightarrow\Tn{i^{\prime}}\twoheadlongrightarrow \Vn{i}\longrightarrow 0,\quad\quad
0\longrightarrow \dVn{i}\hooklongrightarrow\Tn{i^{\prime}}\twoheadlongrightarrow \dVn{i^{\prime}}\longrightarrow 0.
\]
\item[(c)] We have $\Tn i\cong \Tn{i^{\prime}}$ 
iff $i=i^{\prime}$. Moreover, if $M\in\Ob(\T)$ 
is indecomposable, then there 
exists an $i\in\bN$ such that $M\cong \Tn i$.\makeqed
\end{prop}



\begin{proof}
\textbf{(a).} Clear by Proposition~\ref{prop-weyl} and Corollary~\ref{cor-weyl}.


\textbf{(b).} First let $T=\Vn{(a+1)l-1}\otimes \Vn 1$.  
Then $T=\Tn{(a+1)l-1}\otimes \Tn 1$ by (a), and, by~\eqref{eq-tiltfact}, 
we see that $T$ has a $\Delta_q$-filtration giving us the short exact sequence
\[
0\longrightarrow \Vn{(a+1)l}\hooklongrightarrow T\twoheadlongrightarrow \Vn{(a+1)l-2}\longrightarrow 0.
\] 
By duality, $T$ has an analogous $\nabla_q$-filtration, showing that $T$ is tilting. 
It is also indecomposable since otherwise the summand $\Tn{(a+1)l}$ 
of $T$ would be equal to $\Vn{(a+1)l}$, forcing it to be simple, which contradicts 
(b) of Proposition~\ref{prop-weyl}.
This proves the claim in case $b=l-2$.


Now, if $b<l-2$, then consider the diagram
\[
\xy
\xymatrix{
0\ar[r] & \Vn{i^{\prime}}\ar@{^{(}->}[r] & \Vn{i^{\prime}-1}\otimes \Vn 1\ar@{->>}[r]\ar@{^{(}->}[d] & \Vn{i^{\prime}-2}\ar[r] & 0\\
 &  & \Tn{i^{\prime}-1}\otimes \Tn 1\ar@{->>}[d] &  & \\
0\ar[r] & \dVn{i^{\prime}-2}\ar@{^{(}->}[r] & \dVn{i^{\prime}-1}\otimes \dVn 1\ar@{->>}[r] & \dVn{i^{\prime}}\ar[r] & 0.\\
}
\endxy
\]
Here the bottom and top sequences come from~\eqref{eq-tiltfact}. 
A computation shows that these split. This 
gives an inclusion $\Vn{i^{\prime}-2}\hookrightarrow\Tn{i^{\prime}-1}\otimes \Tn 1$ 
which extends to $\Tn{i^{\prime}-2}$, and a surjection 
$\Tn{i^{\prime}-1}\otimes \Tn 1\twoheadrightarrow\dVn{i^{\prime}-2}$ 
factoring through $\Tn{i^{\prime}-2}$. The resulting endomorphism of $\Tn{i^{\prime}-2}$
is non-zero on the 
$i^{\prime}-2$ weight space (the reader unfamiliar 
with this notion might consider Remark~\ref{rem-tria} 
where the notation is $M_i$ for the $i$-th weight space) 
and thus, an isomorphism. 
Hence, 
$\Tn{i^{\prime}-1}\otimes \Tn 1\cong\Tn{i^{\prime}-2}\oplus T$ 
for some $\Uq$-module $T$, which 
is tilting by (a) of Lemma~\ref{lem-tiltcat}. 
Using an inductive argument, we have that $T$ has two $\Delta_q$- 
and $\nabla_q$-factors as in the statement of (b). As above, we see that 
$T=\Tn{i^{\prime}}$, which proves the claim.

\textbf{(c).} To see the first statement note that it suffices to show 
that $\Tn i\not\cong \Tn{i^{\prime}}$ for $i\neq i^{\prime}$. To this end, assume without 
loss of generality that $i^{\prime}>i$. 
By~\eqref{eq-tiltfact2}, we see 
that $\Vn{i^{\prime}}$ appears with multiplicity 
$1$ in $\Tn{i^{\prime}}$, but not in $\Tn i$. Thus,  
$\Tn i\not\cong \Tn{i^{\prime}}$.
That every finite-dimensional 
indecomposable tilting is of this 
form needs a little bit more treatment. 
But it is a standard argument that appears 
in various contexts and can be 
adopted for example 
from~\cite[Section~11.2]{hum} or, alternatively, 
from~\cite[Chapter~II, Section~E.6]{jarag}.
\end{proof}


\begin{ex}\label{ex-tilt1}
Let $l=3$ again. The $\Uq$-tilting module $\Tn 2$ can be 
visualized as
\begin{equation*}
\begin{gathered}
\xymatrix{
 \otimes \ar@{.}[r]\ar@{.}[d]& m_1 \ar@<2pt>[r]^{1}\ar@(ul,ur)|{q^{-1}}  &  m_0  \ar@<2pt>[l]^1\ar@(ul,ur)|{q^{+1}}  \\
 m_1\ar@<2pt>[d]^{1}\ar@(ul,dl)|{q^{-1}} & m_{11} \ar@<2pt>[r]^{1}\ar@<2pt>[d]^{1}\ar@(ul,dl)|{q^{-2}}  &  m_{01}  \ar@<2pt>[l]^1\ar@<2pt>[d]^{1}\ar@{.}[dl]|{+}\ar@(ur,dr)|{q^{0}}  \\
 m_0\ar@<2pt>[u]^{1}\ar@(ul,dl)|{q^{+1}} & m_{10} \ar@<2pt>[r]^{1}\ar@<2pt>[u]^{1}\ar@(ul,dl)|{q^{0}}  &  m_{00}  \ar@<2pt>[l]^1\ar@<2pt>[u]^{1}\ar@(ur,dr)|{q^{+2}}  \\
}
\end{gathered},
\end{equation*}
where we use $m_{ij}=m_i\otimes m_j$. 
By construction, $\Tn 2$ contains $m_{00}$ and it therefore has to be the span of $\{m_{00},q^{-1}m_{10}+m_{01},m_{11}\}$ as indicated above (which is isomorphic to $\Ln 2$).
\end{ex}


\begin{rem}\label{rem-proj}
One can 
show that the category of all 
finite-dimensional projective $\Uq$-modules
is a full 
subcategory of $\T$. 
The same is true for finite-dimensional 
injective $\Uq$-modules: $\T$ has enough injectives and projectives. 
In fact, $\Tn{i}$ is injective and projective 
for all $i\geq l-1$. Hence, 
$\{\Tn{i}\mid i\geq l-1\}$ is a complete set of 
pairwise non-isomorphic, projective-injective
indecomposable, finite-dimensional $\Uq$-modules. 
See for example~\cite[Section~5]{an1}
\end{rem}




\subsection{The linkage principle and blocks}\label{sub-blocks}
Consider the alcove $\mathcal A_0$ 
and its closure $\bar{\mathcal A}_0$ given by
\[
\mathcal A_0=\{k\in\bZ\mid -1<k<l-1\},\quad\quad\bar{\mathcal A}_0=\{k\in\bZ\mid -1\leq k \leq l-1\}.
\]
We call $\mathcal A_0$ the \textit{fundamental alcove}. 
Any other alcove in $\bZ_{\geq -1}$ is 
clearly of the form $\mathcal A_0+il$ for 
some $i\in\bN$. Moreover, we call 
$-1,l-1\in\bar{\mathcal A}_0-\mathcal A_0$ \textit{walls of $\mathcal A_0$}.




The \textit{affine Weyl group} is 
$W_l=\langle s_{r},t_{r}\mid r\in\bZ\rangle$ where $s_{r}$ and $t_r$ act on $\bZ$ via
\[
s_{r}.k=-k-2+(4r+2)l,\quad\quad t_{r}.k=-k-2+4rl,\quad k\in\bZ
\]
Note that $W_l$ is isomorphic 
to the infinite dihedral group $D_{\infty}$: 
we have $W_l\cong \bZ/2\bZ\ltimes l\bZ\cong D_{\infty}$ 
(with $\bZ/2\bZ$ being the Weyl group of $\mathfrak{sl}_2$) and we denote 
the two generators of $D_{\infty}$ by $s$ and $t$ (with the evident association). 
They act ``alcove-wise'' as indicated in~\eqref{eq-alcove}.

 
We denote by $W_l.x$ the orbits 
in $\bZ_{\geq -1}$ under the 
action of the affine Weyl 
group acting on $x\in\bZ_{\geq -1}$. 
For $l=3$ this can be visualized as
\begin{equation}\label{eq-alcove}
\xy
(0,0)*{\includegraphics[scale=0.75]{res/figs/tilt/alcove.eps}};
(-47,0)*{-1};
(-37.5,0)*{0};
(-28,0)*{1};
(-18.5,0)*{2};
(-9.75,0)*{3};
(0,0)*{4};
(9.5,0)*{5};
(18.5,0)*{6};
(28,0)*{7};
(37.5,0)*{8};
(47,0)*{9};
(-18.5,-7.5)*{s\text{-wall}};
(37.5,-7.5)*{s\text{-wall}};
(9.5,7.5)*{t\text{-wall}};
(-18.5,12)*{r=0};
(37.5,12)*{r=1};
(9.5,-12)*{r=1};
(-50,-7.5)*{\text{dead-end}};
\endxy
\end{equation}
with red (or top) action 
by $s$ and green (or bottom) 
action by $t$.



We say $i\in\bN$ is \textit{linked} to $i^{\prime}\in\bN$ 
if there exists $w\in W_l$ such that $w.i=i^{\prime}$. We, by 
convention, set all $\Uq$-modules indexed by negative numbers to be zero.


The following is known as the 
\textit{$\mathfrak{sl}_2$-linkage principle}. 
The more 
general (and highly non-trivial) 
statement can be found in~\cite[Corollaries~4.4 and~4.6]{an2}.



\begin{thm}(\textbf{The linkage principle})\label{thm-link}
All composition factors of $\Tn i$ have maximal weights $i^{\prime}$ linked to 
$i$. Moreover, $\Tn i$ is a simple $\Uq$-module if $i\in W_l.(\bar{\mathcal A}_0-\mathcal{A}_0)$.

In addition, if $i$ is linked to an element of $\mathcal{A}_0$, then 
$\Tn i$ is a simple $\Uq$-module iff $i\in\mathcal{A}_0$.\makeqed
\end{thm}



\begin{proof}
Use Propositions~\ref{prop-weyl} and~\ref{prop-tilt} and Corollary~\ref{cor-weyl}.
\end{proof}




\begin{nota}\label{no-link}
We \textit{always} use $\lambda$ for 
elements in the fundamental alcove 
$\mathcal{A}_0$ and $\mu$ for elements 
on the walls of the fundamental alcove 
$\bar{\mathcal A}_0-\mathcal A_0$. Moreover, set 
\[
\mathcal A_i=\{i^{\prime}\mid il-1<i^{\prime}<(i+1)l-1\}=\mathcal A_0+il,\quad
\bar{\mathcal A}_i=\{i^{\prime}\mid il-1\leq i^{\prime}\leq (i+1)l-1\}=\bar{\mathcal A}_0+il,
\] 
for each $i\in\bN$. Suppose $\lambda\in\mathcal A_0$ 
and $\mu\in\bar{\mathcal A}_0-\mathcal A_0$. We write 
$\lambda_i$ and $\mu_i$ for the unique elements in 
$\mathcal A_i\cap W_l.\lambda$ and $\bar{\mathcal A}_i\cap W_l.\mu$. 
Note that, if $\mu=\mu_0=-1$, then $\mu_1=\mu_2$, $\mu_3=\mu_4$ and so 
forth, while for $\mu=\mu_0=l-1$ 
we have $\mu_0=\mu_1$, $\mu_2=\mu_3$ etc. See also~\eqref{eq-alcove}.
\end{nota}



We define the tilting generator $\Tn{\bbin}$ and 
its cut-off's $\Tn{\leq \bbm}$ that will play a 
crucial role in Section~\ref{sec-quiver}. By 
abuse of notation, we denote them in the same 
way for all $\lambda$ and $\mu$ and hope that 
it is clear from the context which ones we consider.


\begin{defn}(\textbf{The tilting generator})\label{defn-tiltgen}
Let $m\in\bN\cup\{\infty\}$. We call the $\Uq$-modules given by
\[
\Tn{\bbin}=\bigoplus_{i=0}^{\infty}\Tn{\lambda_i}\quad\text{and}\quad\Tn{\leq \bbm}=\bigoplus_{i=0}^m\Tn{\lambda_i}
\] 
the \textit{tilting generator} and its \textit{$m$-th cut-off}, respectively. 
Likewise for $\mu$'s instead of $\lambda$'s.
\end{defn}



Denote by $\T_{\lambda}$ for $\lambda\in\mathcal A_0$ 
the $\lambda$ \textit{block of} $\T$: all indecomposable summands of 
objects of $\T_{\lambda}$ should be of 
the form $\Tn i$ for $i\in W_l.\lambda$. Note that 
blocks are in general not indecomposable as categories.
We, using a similar notation on walls, note the following.



\begin{lem}\label{lem-block}
We have
\[
\T=\bigoplus_{\lambda\in\mathcal A_0}\T_{\lambda}\oplus \T_{-1}\oplus \T_{l-1}
\]
with semisimple categories 
$\T_{-1}$ and $\T_{l-1}$ equivalent 
to the corresponding $\Uv$-module categories. 
We have for all $\lambda\in\mathcal A_0$
(similar for $\T_{-1}$ and $\T_{l-1}$):
\begin{itemize}
\item[(a)] The categories $\T_{\lambda}$ 
are additive, closed under finite direct sums and under duals.
\item[(b)] The categories $\T_{\lambda}$ 
are full Krull-Schmidt subcategories of $\T$. 
Moreover, every indecomposable tilting 
in $\T_{\lambda}$ is of the form $\Tn i$ for some $i\in\bN$ and $i\in W_l.\lambda$.\makeqed
\end{itemize}
\end{lem}


\begin{proof}
This is now only a 
combination of Theorem~\ref{thm-link} and 
Lemma~\ref{lem-tiltcat}. Note that 
$\T_{-1}$ and $\T_{l-1}$ are equivalent 
to categories of $\Uv$-modules by Proposition~\ref{prop-tilt} part (a).
\end{proof}



\begin{nota}\label{no-linkb}
If it is clear from the context 
which $\lambda$ we consider, then we, 
by abuse of notation, write 
$\Tn{\lambda_i}=\Tn{\bbi}$ for short. Similarly for simple and (dual) 
Weyl modules, and on walls.
\end{nota}



\begin{ex}\label{ex-block}
Take $l=3$ again. Then we only have to consider
$\lambda=0,1$ and $\mu=-1,2$.



The indecomposable tiltings in $\T_0$ are 
$\Tn{i}$ for $i=0,4,6,\dots$ as a look 
at~\eqref{eq-alcove} indicates. For $\T_1$ 
they are $\Tn{i}$ for $i=1,3,7,\dots$. The 
two blocks $\T_{-1}$ and $\T_2$ are semisimple 
and consist of direct sums of $\Ln i\cong\Tn{i}$ 
for $i=5,11,\dots$ and for $i=2,8,\dots$ respectively.
\end{ex}



\begin{rem}\label{rem-tria}
As in the usual case for an 
indeterminate $v$, we have a 
triangular decomposition
$\Uq=\Uq^-\Uq^0\Uq^+$, see for example~\cite[Section~1]{apw}. 
There is a character $\chi_i\colon \Uq^0\to\bC$ for any $i\in\bN$ 
(how the character is 
determined by $i$ can be found in~\cite[Lemma~1.1]{apw}).
We note that the Weyl module $\Vn{i}$ 
has the following \textit{universal property}: 
for any $\Uq$-module $M$ there is an isomorphism of vector spaces
\[
\Hom_{\Uq}(\Vn{i},M)\cong\{m\in M_i\mid E^{(j)}m=0\text{ for all }j\in\bN\},
\]
where $M_i=\{m\in M\mid um=\chi_i(u)m,\; u\in\Uq^0\}$ is the $i$ weight space 
of $M$.
This together with Definition~\ref{defn-weyl}, 
Proposition~\ref{prop-weyl} and Corollary~\ref{cor-weyl} imply that
\[
\Hom_{\Uq}(\Vn{i},\Ln{j})\cong\Hom_{\Uq}(\Vn{i},\dVn{j})\cong\delta_{ij}\bC
\]
for all $i,j\in\bN$. In particular,
\begin{equation}\label{eq-SchurLemma}
\Hom_{\Uq}(\Ln{i},\Ln{j})\cong\delta_{ij}\bC,
\end{equation}
i.e. Schur's Lemma holds in our set-up (although $\bC$ is not 
algebraically closed). Note that this is true for general quantum groups over 
arbitrary fields, see~\cite[Corollary~7.4]{apw}.
In addition\footnote{We point out that this fails in general. An explicit counterexample can be found for example in~\cite[Section~5]{ak}.}, we get
\[
\Hom_{\Uq}(\dVn{i},\Vn{i})\cong\bC.
\]
for all $i\in\bN$ (by using 
Proposition~\ref{prop-weyl} and Corollary~\ref{cor-weyl}).
\end{rem}



We call the following maps 
\textit{up} $u^{\lambda}_i$, 
\textit{down} $d^{\lambda}_i$ and 
\textit{loop} $\varepsilon^{\lambda}_i$ 
(\textit{or simply} $u_i,d_i$ and $\varepsilon_i$) respectively.



\begin{prop}\label{prop-link}
There exist up to scalars unique 
$\Uq$-intertwiners $u^{\lambda}_i,d^{\lambda}_i$ with
\begin{gather*}
u^{\lambda}_i\colon \Tn{\lambda_i}\to\Tn{\lambda_{i+1}},\; i=0,1,\dots,\quad\quad d^{\lambda}_i\colon \Tn{\lambda_i}\to\Tn{\lambda_{i-1}},\; i=1,2,\dots,
\\
u^{\lambda}_{i}\circ u^{\lambda}_{i-1}=0=d^{\lambda}_i\circ d^{\lambda}_{i+1},\; i=1,2,\dots,\quad\quad d^{\lambda}_{i+1}\circ u^{\lambda}_i=\varepsilon_i^{\lambda}=u^{\lambda}_{i-1}\circ d^{\lambda}_{i},\; i=1,2,\dots
\\
\varepsilon_i^{\lambda}\circ\varepsilon_i^{\lambda}=0,\; i=1,2,\dots
,\quad\quad d^{\lambda}_1\circ u^{\lambda}_0=0,
\end{gather*}
for $\lambda\in\mathcal A_0$. 
The equation on the bottom right is called 
the \textit{dead-end relation}.\makeqed
\end{prop}



\begin{proof}
First let us assume that we are not in case of the 
dead-end relation (in fact, we leave it to 
the reader to verify this special case), 
that is, the indices $i$ are at least $1$. 
We want to use Proposition~\ref{prop-weyl} 
and Corollary~\ref{cor-weyl}. Moreover, we 
note that the $\Uq$-morphisms below will 
only be unique up to scalars and their 
precise form does not matter. We only 
assume that they are non-zero. In fact, 
by abuse of notation, we always use the 
same symbols, but the maps are of course different in general.




We consider
\[
\xymatrix{
    & \Ln{\lambda_i}\phantom{.} \ar@{^{(}->}[r] & \dVn{\lambda_i}\phantom{.} \ar@{->>}[r]^a & \Ln{\lambda_{i-1}}
                \ar@{->} `r/8pt[d] `/10pt[l] `^dl[ll]|b `^r/3pt[dll] [dll] \\
             & \Ln{\lambda_{i-1}} \ar@{^{(}->}[r]^c & \Vn{\lambda_{i}} \ar@{->>}[r]^d & \Ln{\lambda_{i}}
                \ar@{->} `r/8pt[d] `/10pt[l] `^dl[ll]|e `^r/3pt[dll] [dll] \\
             & \Ln{\lambda_{i}} \ar@{^{(}->}[r]^f & \dVn{\lambda_{i}} \ar@{->>}[r] & \Ln{\lambda_{i-1}}.
    }
\]
Hence, we see that $\Hom_{\Uq}(\dVn{\lambda_i},\Vn{\lambda_i})$ is 
one dimensional: as noted in Remark~\ref{rem-tria}, the 
hom-spaces between Weyl and dual Weyl modules are at 
most $1$-dimensional. The composite $cba$ ensures 
that the dimension is exactly $1$ since it spans the hom-space.




Likewise for $\Hom_{\Uq}(\Vn{\lambda_i},\dVn{\lambda_i})$ by 
using $fed$. Note that the composite $fedcba=0$ due to 
the fact that the middle row is exact. The same holds when the roles 
of $\Vn{\lambda_i}$ and $\dVn{\lambda_i}$ are exchanged. 
Moreover, note that the morphism from $\dVn{\lambda_i}$ to 
$\Vn{\lambda_i}$ uses only $\Ln{\lambda_{i-1}}$ while the 
other way around uses only $\Ln{\lambda_{i}}$. Thus, up 
to scalars, morphisms from $\dVn{\lambda_i}$ to $\Vn{\lambda_i}$ are 
the ``same'' as morphisms from $\Vn{\lambda_{i-1}}$ to $\dVn{\lambda_{i-1}}$.
We can now construct, by using Proposition~\ref{prop-tilt}, 
up $u_i^{\lambda}$ and 
down $d_i^{\lambda}$ 
as a composition of the following maps.
\[
\xymatrix{
    & \Vn{\lambda_i}\phantom{.} \ar@{^{(}->}[r] & \Tn{\lambda_i}\phantom{.} \ar@{->>}[r]^a & \Vn{\lambda_{i-1}}
                \ar@{->} `r/8pt[d] `/10pt[l] `^dl[ll]|b `^r/3pt[dll] [dll] \\
             & \Vn{\lambda_{i-1}} \ar@{^{(}->}[r]^c & \Tn{\lambda_{i-1}} \ar@{->>}[r]^d & \Vn{\lambda_{i-2}}
             \ar@{->} `r/8pt[d] `/10pt[l] `^dl[ll]|e `^r/3pt[dll] [dll] \\
             & \Vn{\lambda_{i-2}} \ar@{^{(}->}[r]^f & \Tn{\lambda_{i-2}} \ar@{->>}[r] & \Vn{\lambda_{i-3}}.
    }
\]
We define $d_i^{\lambda}$ to be the 
composite $cba$. Similar for $u_i^{\lambda}$, 
but using the right-hand side of part (b) of 
Proposition~\ref{prop-tilt}. By the same 
reasoning as above we see that they are unique up to scalars.




We have to check the relations between the 
various up $u_i^{\lambda}$ and down $d_i^{\lambda}$ maps 
now. To see that $u^{\lambda}_{i}\circ u^{\lambda}_{i-1}=0$ 
and $d^{\lambda}_i\circ d^{\lambda}_{i+1}=0$ we can simply 
use the second diagram above and its dual 
counterpart and the fact that the rows are exact. Now consider
\[
\xymatrix{
\Tn{\lambda_{i-1}}\ar@{<->}[d]|{\mathrm{id}} & \Vn{\lambda_{i-1}}\ar@{_{(}->}[l]\ar[d]|{\Ln{\lambda_{i-1}}\to\Ln{\lambda_{i-1}}} & \Tn{\lambda_{i}}\ar@{<->}[d]|{\mathrm{id}}\ar@{->>}[r]\ar@{->>}[l] & \dVn{\lambda_{i}}\ar@{^{(}->}[r]\ar[d]|{\Ln{\lambda_{i-1}}\to\Ln{\lambda_{i-1}}} & \Tn{\lambda_{i+1}}\ar@{<->}[d]|{\mathrm{id}}\\
\Tn{\lambda_{i-1}}\ar@{->>}[r] & \dVn{\lambda_{i-1}}\ar@{^{(}->}[r] & \Tn{\lambda_{i}} & \Vn{\lambda_{i}}\ar@{_{(}->}[l] & \Tn{\lambda_{i+1}}\ar@{->>}[l].
}
\]
Combining everything, we see that 
(up to scalars) $d^{\lambda}_{i+1}\circ u^{\lambda}_i=\varepsilon_i^{\lambda}=u^{\lambda}_{i-1}\circ d^{\lambda}_{i}\neq 0$. Moreover, by 
the reasoning above, the $\varepsilon_i^{\lambda}$ 
is an up to scalars unique non-zero 
$\Uq$-morphism $\Tn{\lambda_i}$ to $\Tn{\lambda_i}$ that squares to zero.
\end{proof}

\newpage

\begin{cor}\label{cor-homspaces}
Let $i,i^{\prime}\in\bN$. Then we have the following.
\begin{itemize}
\item[(a)] Outside of walls:
\[
\Hom_{\Uq}(\Tn{\lambda_i},\Tn{\lambda_{i^{\prime}}})\cong\begin{cases}\bC[\varepsilon],  & \text{if }|i-i^{\prime}|=0\text{ and }i=i^{\prime}\neq 0,\\
  \bC,  & \text{if }|i-i^{\prime}|=1 \text{ or }i=i^{\prime}=0,\\
  0,  & \text{if }|i-i^{\prime}|>1,\end{cases}
\]
where $\bC[\varepsilon]\cong \bC[X]/(X^2)$ denotes the $\bC$-algebra of dual numbers.
\item[(b)] On walls:
\[
\Hom_{\Uq}(\Tn{\mu_i},\Tn{\mu_{i^{\prime}}})\cong\begin{cases}\bC,  & \text{if }|i-i^{\prime}|=0,\\
  0,  & \text{if }|i-i^{\prime}|>0,\end{cases}
\]
where $\bC$ is the trivial $\Uq$-module.\makeqed
\end{itemize}
\end{cor}



\begin{proof}
Because of Proposition~\ref{prop-link}, we only 
need to verify (b). But since the $\Tn{\mu_i}$ 
are simple on walls (by Proposition~\ref{prop-tilt}), 
part (b) follows from ``Schur's Lemma''~\eqref{eq-SchurLemma}.
\end{proof}



\subsection{Translation functors}\label{sub-translation}
Fix $\lambda\in\mathcal A_0$ and 
$\mu\in\bar{\mathcal{A}}_0-\mathcal A_0$. Set $\nu=|\lambda-\mu|$ 
(in fact, $\nu=|\lambda_i-\mu_i|$ for all 
$i\in\bN$ and $\nu$ will stay in $\bar{\mathcal{A}}_0$). 
Recall that there is a unique simple 
$\Uq$-module $\Ln{\nu}\cong\Tn{\nu}$.
Denote by $p_{\lambda}\colon\T\to\T_{\lambda}$ the 
\textit{projection onto the block} $\T_{\lambda}$ functor. Similarly for $\mu$.



\begin{defn}(\textbf{Onto and out of the wall})\label{defn-onout}
Given $\lambda$ and $\mu$, we define 
two functors via
\[
\mathcal T_{\lambda}^{\mu}\colon\T_{\lambda}\to\T_{\mu},\;M\mapsto p_{\mu}(M\otimes_{\bC}\Tn{\nu}),\quad\quad\mathcal T^{\lambda}_{\mu}\colon\T_{\mu}\to\T_{\lambda},\;M\mapsto p_{\lambda}(M\otimes_{\bC}\Tn{\nu}).
\]
We call them \textit{translation onto} the $\mu$-wall $\mathcal T_{\lambda}^{\mu}$ and 
\textit{translation out of} the $\mu$-wall $\mathcal T^{\lambda}_{\mu}$.
\end{defn}




\begin{defn}(\textbf{Through the wall})\label{defn-trans}
Define $\Theta_s^{\lambda}=\mathcal T^{\lambda}_{-1}\circ \mathcal T_{\lambda}^{-1}$ and $\Theta_t^{\lambda}=\mathcal T^{\lambda}_{l-1}\circ \mathcal T_{\lambda}^{l-1}$. We 
call them \textit{translation through the wall functors}.
\end{defn}



If it is clear which $\lambda$ we are using we, abusing notation, denote them by $\Theta_s,\Theta_t$.



\begin{prop}\label{prop-functors}
For all $i\in\bN$ we have 
the following.
\begin{itemize}
\item[(a)] The functors 
$\mathcal T_{\lambda}^{\mu}$ and $\mathcal T^{\lambda}_{\mu}$ are well-defined (their definition gives $\Uq$-tilting modules in the right blocks), 
adjoints (left and right) and 
exact. Thus, $\Theta^{\lambda}_s$ and 
$\Theta^{\lambda}_t$ are exact and self-adjoint.
\item[(b)] We have (recalling that $\Ln{\mu_i}\cong\Tn{\mu_i}$)
\[
\mathcal T_{\lambda}^{\mu}(\Tn{\lambda_i})\cong\begin{cases}\Tn{\mu_{i-1}}\oplus\Tn{\mu_{i+1}},  & \text{if }\mu_i>\lambda_i,\\
  \Tn{\mu_i}\oplus\Tn{\mu_i},  & \text{if }\mu_i<\lambda_i,\end{cases}\quad
\mathcal T^{\lambda}_{\mu}(\Tn{\mu_i})\cong\begin{cases}\Tn{\lambda_{i+1}},  & \text{if }\mu_i>\lambda_i,\\
 \Tn{\lambda_{i}},  & \text{if }\mu_i<\lambda_i.\end{cases}
\]
\item[(c)] The dead-end relations $\Theta^{\lambda}_s(\Tn{\lambda_0})\cong0$, $\Theta^{\lambda}_s(\Tn{\lambda_1})\cong\Tn{\lambda_2}$, and $\Theta^{\lambda}_t(\Tn{\lambda_0})\cong\Tn{\lambda_1}$. Moreover, we have
\[
\Theta^{\lambda}_{s\text{ or }t}(\Tn{\lambda_i})\cong\begin{cases}\Tn{\lambda_{i-1}}\oplus\Tn{\lambda_{i+1}} ,  & \text{if }i>1\text{ is odd for }s\text{ and even for }t,\\
  \Tn{\lambda_{i}}\oplus\Tn{\lambda_{i}} ,  & \text{if }i>0\text{ is odd for }t\text{ and even for }s.\end{cases}
\]
\end{itemize}
Here we set $\Tn{-1}=\Tn{\lambda_{-1}}=\Tn{\mu_{-1}}=0$.\makeqed
\end{prop}


Note that biadjoint functors have in 
addition some other nice properties, see e.g.~\cite[Section~2]{kh5}. Moreover, 
in the 
case $\mu=-1$ we have $\mu_i>\lambda_i$ 
iff $i$ is odd whereas in the case $\mu=l-1$ 
we have $\mu_i>\lambda_i$ iff $i$ is even as 
a look at~\eqref{eq-alcove} should convince the reader.



\begin{proof}
\textbf{(a).} The functors are well-defined by 
(c) of Proposition~\ref{prop-basicp} and
part (b) of Lemma~\ref{lem-tiltcat}, i.e. tensor products of tiltings are tiltings. 
The other statements can be verified as in the 
case of category $\mathcal O$, i.e. we can adopt~\cite[Sections~7.1 and~7.2]{hum} without difficulties. 




\textbf{(b).} Use 
Propositions~\ref{prop-weyl} 
and~\ref{prop-tilt}, Corollary~\ref{cor-weyl} and Theorem~\ref{thm-link}.



\textbf{(c).} This is just a direct application of the finer list of 
statements given in (b).
\end{proof}



We use the convention that 
$\Theta^{\lambda}_{k-tst}$ denotes an 
\textit{alternating} composition of length 
$k$ of translation through the wall functors starting with 
$\Theta^{\lambda}_{t}$. Likewise for $\Theta^{\lambda}_{k-sts}$.
Using (c) of Proposition~\ref{prop-functors}, we get the following.


\begin{cor}(\textbf{Combinatorics of the translation 
through the wall functors})\label{cor-theta}
We have:
\begin{itemize}
\item[(a)] $\Theta_s^{\lambda}\circ\Theta_s^{\lambda}\cong \Theta_s^{\lambda}\oplus \Theta_s^{\lambda}$ and $\Theta_t^{\lambda}\circ\Theta_t^{\lambda}\cong \Theta_t^{\lambda}\oplus \Theta_t^{\lambda}$ as functors.
\item[(b)] We have for $i\in\bN$ even respectively odd, 
that (with $k\geq 0$ terms $\Theta^{\lambda}_{t\text{ or }s}$) 
there exist multiplicities $m_j\in\bN$ (that can be zero) such that
\begin{gather*}
\Theta^{\lambda}_{k-tst}\Tn{\lambda_i}=(\Theta^{\lambda}_{s\text{ or }t}\circ\dots\Theta^{\lambda}_t\circ\Theta^{\lambda}_s\circ\Theta^{\lambda}_t)\Tn{\lambda_i}\cong \Tn{\lambda_{i+k}}\oplus \bigoplus_{j<i+k}\Tn{\lambda_{j}}^{\oplus m_j},
\\
\Theta^{\lambda}_{k-tst}\Tn{\lambda_i}=(\Theta^{\lambda}_{s\text{ or }t}\circ\dots\Theta^{\lambda}_t\circ\Theta^{\lambda}_s\circ\Theta^{\lambda}_t)\Tn{\lambda_i}\cong \Tn{\lambda_{i-1+k}}^{\oplus 2}\oplus \bigoplus_{j<i-1+k}\Tn{\lambda_{j}}^{\oplus m_j}.
\end{gather*}
Similarly for $i\in\bN$ odd respectively even and $\Theta^{\lambda}_{k-sts}$.\makeqed
\end{itemize}
\end{cor}

\begin{proof}
Clear from Proposition~\ref{prop-functors}.
\end{proof}

Define the following functors:
\begin{gather*}
\begin{aligned}
\mathcal T^{\mu,i}_{\lambda}=p_{\bbi}\circ\mathcal T^{\mu}_{\lambda},
\quad\quad&
\mathcal T_{\mu}^{\lambda,i}(\Tn{\mu_j})=\begin{cases}
\delta_{i,j}\Tn{\lambda_{i}},&\text{if }\mu=-1, i\text{ even},\\
\delta_{i,j}\Tn{\lambda_{i}},&\text{if }\mu=l-1, i\text{ odd},
\end{cases}
\\
\Theta_{s\text{ or }t}^{\lambda,i}&=\begin{cases}p_{\bbi}\circ\Theta^{\lambda}_{s}, &\text{if }i\text{ is even},\\
p_{\bbi}\circ\Theta^{\lambda}_{t}, &\text{if }i\text{ is odd}.
\end{cases}
\end{aligned}
\end{gather*}
Here $p_{\bbi}$ denotes the projection to the $\Tn{\bbi}$-part.

\begin{defn}\label{defn-tiltcatendo}
Given a $\Uq$-tilting module $T$, denote by 
$\mathcal F_{T}$ any functor $\mathcal F_{T}\cong \cdot\otimes_{\bC}T$.
A functor $\mathcal F\colon\T\to\T$ is called 
\textit{projective} if it is a direct summand 
of functors of the form $\mathcal F_{T}$.
We denote by $\Endo(\T)$ the \textit{category 
of projective endofunctors} $\mathcal F\colon\T\to\T$ 
whose morphisms $\Hom_{\Endo(\T)}(\mathcal F,\mathcal F^{\prime})$ 
are natural transformations $\eta\colon\mathcal F\to\mathcal F^{\prime}$. 
Moreover, we denote by $\Endo(\T_{\lambda})$ the \textit{category 
of projective endofunctors} $\mathcal F\colon\T_{\lambda}\to\T_{\lambda}$ 
that can be obtained via compositions and countable direct sums 
of functors of the form $\Theta_{s\text{ or }t}^{\lambda,i}$. Similarly for walls $\mu$. 
Here we count the empty sum as the identity functor. 
\end{defn}


It follows from Lemma~\ref{lem-profunc} below that these are actually categories. 
Moreover, by construction, the translation through the wall functors 
$\Theta_s^{\lambda}$ or $\Theta_t^{\lambda}$ are objects in $\Endo(\T_{\lambda})$, 
since
\begin{equation}\label{eq-projectivefunc}
\Theta_s^{\lambda}=\!\!\bigoplus_{i\in\bZ_{\geq 1}}\!\!\Theta_s^{\lambda,2i},\quad\quad\Theta_t^{\lambda}=\!\!\bigoplus_{i\in\bZ_{\geq 0}}\!\!\Theta_s^{\lambda,2i+1},
\end{equation}
as follows by Proposition~\ref{prop-functors}. Similarly,
\begin{equation}\label{eq-projectivefunc2}
\mathcal T_{\lambda}^{\mu}=\!\!\bigoplus_{i\in\bZ_{\geq 0}}\!\!\mathcal T_{\lambda}^{\mu,i},\quad\quad \mathcal T_{\mu}^{\lambda}=\begin{cases}\bigoplus_{i\in\bZ_{\geq 0}}\!\!\mathcal T_{\mu}^{\lambda,2i}, &\text{if }\mu=-1,\\
\bigoplus_{i\in\bZ_{\geq 0}}\!\!\mathcal T_{\mu}^{\lambda,2i+1}, &\text{if }\mu=l-1.\end{cases}
\end{equation}

\newpage

\begin{lem}\label{lem-profunc}
$\Endo(\T), \Endo(\T_{\lambda}), \Endo(\T_{\mu})$ are additive 
categories of exact functors.\makeqed
\end{lem}



\begin{proof}
That $\Endo(\T)$ is closed under composition 
follows from part (b) of Lemma~\ref{lem-tiltcat}, 
and the identity functor is projective because 
of $\mathrm{id}\cong\cdot\otimes_{\bC} \Tn 0$. Hence, $\Endo(\T)$ is a category.
That $\Endo(\T)$ is preserved under finite 
direct sums (showing that the category is 
additive) follows from part (a) of Lemma~\ref{lem-tiltcat}. 
The statement for the $\lambda$ and $\mu$ 
versions follow by construction.
Last, that functors in 
$\Endo(\T)$ are exact follows because tensoring over a 
field always respects exact sequences. 
Likewise for the $\lambda$ and $\mu$ versions which come 
from projections of such functors.
\end{proof}



The following lemma 
is true in more generality (see for example~\cite[Chapter~4, Section~6]{ass}), 
but we restrict to our case here. Recall that a 
functor between additive categories 
$\mathcal F\colon\mathcal C\to\mathcal D$ is 
called \textit{indecomposable} if any 
decomposition $\mathcal F\cong\mathcal F_1\oplus\mathcal F_2$ 
implies that $\mathcal F_1\cong 0$ or $\mathcal F_2\cong 0$.
Moreover, we call a category \textit{infinite} Krull-Schmidt 
if each object can be uniquely decomposed (up to permutation) 
into a countably infinite direct sum of indecomposable objects.



\begin{lem}\label{lem-profuncinde}
A functor $\mathcal F\cong \cdot\otimes_{\bC} T$ is 
indecomposable iff $T$ is an indecomposable 
$\Uq$-module. Thus, $\Endo(\T)$ is Krull-Schmidt. In contrast, the $\lambda$ and 
$\mu$ versions are infinite Krull-Schmidt.\makeqed
\end{lem}



\begin{proof}
Assume that $T$ decomposes 
into $T_1\oplus T_2$. Then $\mathcal{F}\cong (\cdot\otimes_{\bC}T_1)\oplus (\cdot\otimes_{\bC}T_2)$.
On the other hand, by Yoneda 
and the tensor-hom adjunction, 
the $T$ representing $\mathcal F$ is 
uniquely determined up to isomorphism. 
Thus, a decomposition of $\mathcal F$ 
induces a decomposition of $T$. This 
implies, by using 
part (a) of Lemma~\ref{lem-tiltcat}, that $\Endo(\T)$ is Krull-Schmidt.

The statement for the $\lambda$ and $\mu$ 
versions follow again by construction.
\end{proof}





\begin{prop}\label{prop-profunc}
The functors $\mathcal T^{\mu,i}_{\lambda}$, $\mathcal T_{\mu}^{\lambda,i}$ and 
$\Theta^{\lambda,i}_{s\text{ or }t}$ are 
all indecomposable.
\end{prop}



\begin{proof}
Any non-trivial decomposition 
$\Theta^{\lambda,i}_{s\text{ or }t}=F_1\oplus F_2$ 
gives rise to a non-trivial idempotent 
$\eta\in\Endo(\Theta^{\lambda,i}_{s\text{ or }t})$. 
By using Proposition~\ref{prop-functors} together 
with Corollary~\ref{cor-homspaces} we see that 
such a natural transformation can not exist 
(most of the $\Uq$-intertwiners are nilpotent). 
Let us illustrate this in an explicit calculation for $\Theta^{\lambda,i}_t$ 
with $i\neq 0$.
Applying $\Theta^{\lambda,i}_t$ on the sequence
\begin{equation}\label{eq-inde}
\xy
\xymatrix{
 \Tn{\bbid}
\ar@<2pt>[r]^/0.25cm/{u_{i-1}} & \Tn{\bbi}\ar@<2pt>[l]^/-0.25cm/{d_{i}}\ar@<2pt>[r]^/-0.25cm/{u_i} 
& \Tn{\bbiu}\ar@<2pt>[l]^/0.25cm/{d_{i+1}}
}
\endxy
\end{equation}
gives us (with bottom and top row given by applying 
$\Theta^{\lambda,i}_t$ to~\eqref{eq-inde})
\begin{gather}\label{eq-inde2}
\begin{aligned}
\xy
\xymatrix{
\Tn{\bbi}\ar@<2pt>[r]^/-1.0em/{\begin{pmatrix} 0\\ 1 \end{pmatrix}}
\ar@{<->}[dd]|{\eta_{\Tn{\bbid}}}
&
\Tn{\bbi}\oplus\Tn{\bbi}\ar@<2pt>[r]^/0.9em/{\begin{pmatrix} 1 & 0\end{pmatrix}}
\ar@<2pt>[l]^/0.9em/{\begin{pmatrix} 1 & 0 \end{pmatrix}}
\ar@{<->}[dd]|{\eta_{\Tn{\bbi}}}
&
\Tn{\bbi}\ar@<2pt>[l]^/-0.9em/{\begin{pmatrix} 0 \\ 1 \end{pmatrix}}
\ar@{<->}[dd]|{\eta_{\Tn{\bbiu}}}\\
& & &\\
\Tn{\bbi}\ar@<2pt>[r]^/-1.0em/{\begin{pmatrix} 0\\ 1 \end{pmatrix}}
&
\Tn{\bbi}\oplus\Tn{\bbi}\ar@<2pt>[r]^/0.9em/{\begin{pmatrix} 1 & 0\end{pmatrix}}
\ar@<2pt>[l]^/0.9em/{\begin{pmatrix} 1 & 0 \end{pmatrix}}
&
\Tn{\bbi}.\ar@<2pt>[l]^/-0.9em/{\begin{pmatrix} 0 \\ 1 \end{pmatrix}}
}
\endxy
\end{aligned}
\end{gather}
Here we used Proposition~\ref{prop-functors} to calculate 
$\Theta^{\lambda,i}_t(\Tn{i})$, while the images of the 
$\Uq$-intertwiners can be calculated as in Example~\ref{ex-howtofunctor}. 
Note that $\eta$ above is assumed to be a natural transformation, i.e. 
all squares in~\eqref{eq-inde2} commute.
Now, under the assumption that $\eta$ is an idempotent, we have $\eta_{\Tn{\bbid}}$ 
equals $0$ or $1$ because of Corollary~\ref{cor-homspaces}. 
In the first case we obtain that all others components 
of $\eta$ are zero as well, while in the second case we obtain 
that all others components 
of $\eta$ are the identity as well 
(both follow via direct computation using Corollary~\ref{cor-homspaces}, the assumption 
that all components of $\eta$ are idempotents 
and the commutativity of all squares in~\eqref{eq-inde2}). 
The same arguments and calculations show the 
claim for the functors $\Theta^{\lambda,i}_s$ 
and $\mathcal T_{\lambda}^{\mu,i}$ 
(except that our job is even easier for the latter since all 
involved tiltings are simple). 
The claim for the functor $\mathcal T^{\lambda,i}_{\mu}$ is left to the reader.
This shows 
the statement.
\end{proof}

\begin{rem}\label{rem-projective}
The category $\Endo(\T_{\lambda})$ is the Karoubi envelope 
(in the sense of Definition~\ref{defn-karoubi}) of the 
category $\Endo_{\Theta}(\T_{\lambda})$ generated 
(via direct sums and compositions) by the 
translation through the wall functors $\Theta^{\lambda}_{s\text{ or }t}$. 
This follows from~\eqref{eq-projectivefunc} and Proposition~\ref{prop-profunc}. 
We call $\Endo(\T_{\lambda})$ the category of projective endofunctors 
although it would be historical more accurate to call 
$\Endo_{\Theta}(\T_{\lambda})$ the category of projective endofunctors. 
This convention is due to the fact that, in contrast to the case of category $\mathcal O$ 
(see e.g.~\cite{bg}), our buildings blocks are the 
various $\Theta^{\lambda,i}_{s\text{ or }t}$ 
and not the functors $\Theta^{\lambda}_{s\text{ or }t}$ themselves.
\end{rem}